\documentclass[quick,con]{debate}
\motion{AI的迅猛发展提升 / 降低了人类创作者存在的意义}
\stance{反方}{降低了}
\begin{document}
\debatetitle
\begin{thesisbox}AI 让社会对多数创作者的需要绕开了人，留下来的人手里只剩修补：用 AI，作品里的你变少；不用 AI，需要你的人变少。\end{thesisbox}

\section{定义与判准（标准表述）}
\begin{fields}
\field{AI 的迅猛发展}{2022 年以来，机器生成文字、图像、音乐和声音的能力突然变强，而且可以不经过人直接交出成品。}
\field{人类创作者}{把创造文艺作品当成自己一部分的人，职业的、业余的都算；尤其要看不出名、靠作品被需要的大多数。}
\field{存在的意义}{作用与重要性（汉典）。两层：这个世界还有多需要、多珍视由人来创作；创作还能给创作者自己留下多少。}
\field{判准}{和 AI 爆发之前比，由人来创作这件事整体变重还是变轻；需要和珍视两半都算，看多数创作者。\sub{我方要证明}{需要在多数人身上缩水，珍视只落在少数；作品里属于创作者的部分在变少。}\sub{对方要证明}{珍视和需要加在一起，对多数创作者变重了。}}
\field{承认线}{法律和标注确实把人和机器分开了，门槛也确实降了。但被区分不等于被需要。}
\field{偏向定义 30 秒口径}{把意义定义成“只有人能给的东西”，等于说 AI 替代什么都不算损失，是把剩下的当成全部。词典里的意义是作用与重要性，被需要也算。}
\end{fields}

\section{我方论点}
\begin{tabularx}{\linewidth}{@{}L{5em}Y Y L{6.4em}@{}}\toprule
\thead{标签}&\thead{一句话}&\thead{核心数据 / 学理 / 案例}&\thead{出处}\\\midrule
\textbf{不再被需要}&对不出名的大多数，被需要是意义最实的一块；这份需要绕开了人，连最好的也挡不住&数据：写作类月收入 -5.2\%，图像类 -9.4\%，顶尖者受冲击更大；26\% 插画师、36\% 译者已丢活。学理：溢价只流向少数。案例：配音师声音被做成产品，判赔 25 万元&Hui 等 2024；英国作家协会 2024；Rosen 1981；北京互联网法院 2024\\
\textbf{手艺被掏空}&留下来的人，作品里属于自己的部分在变少；构思去了甲方，执行给了机器&数据：平均内容新颖度下降，视觉新颖度全面下降；用 AI 的故事更趋同。学理：构思与执行分离。案例：Amber Yu 修图，价钱是原来十分之一&Zhou \& Lee 2024；Doshi \& Hauser 2024；Braverman 1974；Rest of World 2023\\
\bottomrule\end{tabularx}

\section{对方论点 → 一句话反驳}
\begin{tabularx}{\linewidth}{@{}Y Y Y@{}}\toprule
\thead{对方会说}&\thead{我方一句话回}&\thead{对方追问时}\\\midrule
人味成了稀缺品，社会在选人&被区分不等于被需要；标签保护的是不被冒名，不是有人来请&“76\% 想分辨不是在乎？”→ 在乎别被骗，和请你画是两回事\\
立意回到人&立意回到的是甲方，不是画师&“甲方也是创作者”→ 那就承认把创作当志业的人被替换了\\
法律只认人是作者&机器不能署名，不等于有人来请你&“保护就是珍视”→ 一道防护栏，不是一块奖牌\\
门槛降低，更多人能创作&人人都能生成，创作者这个身份越来越轻&“表达变容易不好吗”→ 好，但不等于以创作为志业的人更重要\\
订单只降了 2\%&收入降 5.2\%，图像类 9.4\%，顶尖者受冲击更大&“那是短期”→ 方向清楚，两成六的插画师已经丢了活\\
\bottomrule\end{tabularx}

\section{必问与必答}
\begin{points}{必问（对方不答就复述一次，不换题）}
\item 正方说的被珍视，有多少落到了一个不出名的插画师身上？（全场追问）
\item 被珍视却没人来请，这算意义提升吗？
\item 修光线、修手指，算不算立意回到画师？
\end{points}
\begin{points}{必答（15 秒正面回答）}
\item 贵方说的意义包括被珍视吗？ → 包括，两半都算；珍视落在少数，需要在多数身上缩水，净额变轻。
\item 八成的人要求标注，不是在乎人吗？ → 在乎的是别被骗；在乎，和请你画，是两回事。
\item 到 2028 年那组数字是预测吧？ → 是预测，我方标明；已经发生的看平台研究和作者调查。
\end{points}

\section{对辩战场概要（三场）}
\begin{tabularx}{\linewidth}{@{}L{5.4em}Y Y Y@{}}\toprule
\thead{对辩}&\thead{开场问题}&\thead{目标}&\thead{收束语}\\\midrule
反三（第一战场）&珍视落在谁身上？&把“两半都算”打成“珍视只给少数”&需要在多数人身上缩水\\
反一（中段收拢）&给我一个不出名的人更被需要的证据&拉回看多数，逼正方交证据&正方讲的是规则和少数\\
反二（结辩前）&用 AI 你变少，不用 AI 需要你的人变少，正方站哪头？&把全场收成夹击&两头都在变轻\\
\bottomrule\end{tabularx}

\section{如果……就……}
\begin{contingency}
\ifthen{对方把意义定义成“只有人能给的”}{背口径：那是把剩下的当成全部；被需要也是作用与重要性。}
\ifthen{对方拿出没预料到的数据}{先问出处和样本，再问它落在多数还是少数、是被需要还是只是被区分。}
\ifthen{“手艺被掏空”被打穿}{收缩到“不再被需要”，用“顶尖者受冲击更大”说明质量换不来被需要。}
\end{contingency}

各环节：质询全是 90 秒双边，不频繁打断；反二质询正一为反一开头取材；反二申论是全场第一篇；反三申论先回应正三；反一先结辩，封路不猜台词。

\begin{recordcard}
\record{反二质询结束}{质询拿到了什么}{\opt{质询拿到了正方承认「执行型订单在减少」}\opt{质询里正方没有正面回答「被珍视却没人来请，算不算提升」}\opt{质询没有拿到明确成果}}{反一立论·开头}
\record{正一立论后}{正方立论的主干}{\opt{正方立论建立在「人味稀缺，所以更被珍视」上}\opt{正方立论建立在「立意回到人」上}\opt{正方立论建立在「法律和标注只认人」上}\opt{正方立论的主干不在以上几条时}}{反二申论·拆正方立论}
\record{正三申论后}{正三申论主打什么}{\opt{正三申论主打「订单少了不等于意义轻了」}\opt{正三申论主打「修图的是手，立意还在人」}\opt{正三申论主打「更多人能创作」}\opt{正三申论的重点不在以上几条时}}{反三申论·回应正三}
\record{正三申论后}{正方到此最有力的一点}{\opt{正方全场最有力的是「法律和标注只认人」}\opt{正方全场最有力的是「贴上人类标签评价更高」}\opt{正方全场最有力的是「AI 让普通人也能创作」}\opt{正方最有力的一点不在以上几条时}}{反三申论·处理正方最强点}
\record{二辩对辩后}{全场主战场落在}{\opt{全场主战场落在「需要和珍视哪个算数」}\opt{全场主战场落在「修图工还是立意者」}\opt{主战场不清楚时}}{反一结辩·归纳分歧}
\record{二辩对辩后}{正方全场最有力的一点}{\opt{正方全场最有力的是「法律和标注只认人」}\opt{正方全场最有力的是「AI 让普通人也能创作」}\opt{正方最有力的一点不在以上几条时}}{反一结辩·处理最强点}
\record{二辩对辩后}{正方全场主打的路线}{\opt{正方全场主打「人味稀缺所以更珍贵」}\opt{正方全场主打「意义不等于订单」}\opt{正方全场主打「一个重新能创作的人」}\opt{正方全场路线不在以上几条时}}{反一结辩·封路}
\record{全场}{对方回避了哪些问题}{（写下来）}{申论与结辩里点名}
\end{recordcard}
\end{document}
